%Root main.tex
%---------------------------------------------------------------------------
%付録B　付録
%----------------------------------------------------------------------------
\chapter{生成したHDLコードの仕様}
\thispagestyle{fancy}
\label{chap:コード}
\section{MSDB\_Control}
トップモジュール(Simulink全体)であり，入出力の信号名、ビット数などを定義している．入出力を表\ref{for:MSDB_Control}に示す．
\begin{table}[h]
 \center
 \caption{MSDB\_Controlポート一覧}
 \label{for:MSDB_Control}
 \small
 \begin{tabular}{c|c|c|c} \hline \hline
  信号名            &   方向    &   bit幅   &   説明   \\ \hline
  Omega\_Rpm\_Digital &   In      &   16      &   モータ回転数   \\ \hline
  Count\_2500       &   In      &   16      &   カウンタ信号   \\ \hline
  I\_d              &   In      &   16      &   d軸実電流   \\ \hline
  I\_q              &   In      &   16      &   q軸実電流   \\ \hline
  I\_dref           &   In      &   16      &   d軸指令電流   \\ \hline
  I\_qref           &   In      &   16      &   q軸指令電流   \\ \hline
  Vd                &   Out     &   16      &   d軸操作量   \\ \hline
  Vq                &   Out     &   16      &   q軸操作量   \\ \hline
 \end{tabular}
\end{table}

%\inputminted{vhdl}{image/appendixC/MSDB_Control.vhd}
\begin{center}
\label{lisC.1}
    ソースコードC.1：DC-MSSRDB\_Control
\lstinputlisting[language=VHDL]{image/appendixC/MSDB_Control.vhd}
\end{center}
%%%%%%%%%%%%%%%%%%%%%%%%%%%%%%%%%%%%%%%%%%%%%%%%%%%%%%%%%%%%%%%%%%%%%%%%%%%%%%%%%%%%%%